\chapter{Similarity Transformations and Jordan Canonical Forms}

\section{Dual Representation of State-Space Models}
Before exploring state transformations, it is important to note that the state-space representation of a system is not unique. A particularly important alternative representation is the \textbf{dual representation} (or transpose representation).

Consider a state-space model:
\begin{align}
    \dot{x}(t) &= A x(t) + B u(t) \\
    y(t) &= C x(t) + D u(t)
\end{align}
Its dual state-space representation is defined as:
\begin{align}
    \dot{\tilde{x}}(t) &= A^\trp \tilde{x}(t) + C^\trp u(t) \\
    y(t) &= B^\trp \tilde{x}(t) + D u(t)
\end{align}
The dual system swaps the role of input and output matrices ($B \leftrightarrow C^\trp$) and transposes the state transition dynamics. As we will see, duality is a mathematical cornerstone when linking controllability with observability, and controller design with observer design.

\section{Similarity Transformations}
A similarity transformation changes the coordinate system of the state space without altering the input-output behavior of the system. Let $T \in \R^{n \times n}$ be a non-singular (invertible) matrix, and define a new state vector $\tilde{x}(t)$ by the relation:
\begin{equation}
    x(t) = T \tilde{x}(t) \iff \tilde{x}(t) = T^{-1} x(t)
\end{equation}
Substituting this relation into the original LTI state-space equations yields:
\begin{align}
    T \dot{\tilde{x}}(t) &= A T \tilde{x}(t) + B u(t) \nonumber \\
    \dot{\tilde{x}}(t) &= T^{-1} A T \tilde{x}(t) + T^{-1} B u(t) \\
    y(t) &= C T \tilde{x}(t) + D u(t)
\end{align}
Thus, the system in the new coordinates is represented by the similar matrices:
\begin{equation}
    \tilde{A} = T^{-1} A T, \quad \tilde{B} = T^{-1} B, \quad \tilde{C} = C T, \quad \tilde{D} = D
\end{equation}

\subsection{Properties of Similarity Transformations}
Similarity transformations preserve the fundamental physical and mathematical properties of the system:
\begin{mytheorem}{Invariance of Key Operators under Similarity}
For any non-singular matrix $T \in \R^{n \times n}$ and similar matrix $\tilde{A} = T^{-1} A T$, the following invariants hold:
\begin{enumerate}
    \item \textbf{Eigenvalues}: The eigenvalues of $\tilde{A}$ are identical to those of $A$, i.e., $\lambda\{\tilde{A}\} = \lambda\{A\}$.
    \item \textbf{Characteristic Polynomial}: $\det(sI - \tilde{A}) = \det(sI - A)$.
    \item \textbf{Transfer Function (Input-Output Invariance)}:
    \begin{equation}
        \tilde{C}(sI - \tilde{A})^{-1}\tilde{B} + \tilde{D} = C(sI - A)^{-1}B + D
    \end{equation}
\end{enumerate}
\end{mytheorem}

\begin{proof}
We provide the step-by-step mathematical proofs:
\begin{enumerate}
    \item \textbf{Characteristic Polynomial \& Eigenvalues}:
    \begin{align}
        \det(sI - \tilde{A}) &= \det(sI - T^{-1}AT) = \det(sT^{-1}IT - T^{-1}AT) \nonumber \\
        &= \det(T^{-1}(sI - A)T) = \det(T^{-1}) \det(sI - A) \det(T) \nonumber \\
        &= \frac{1}{\det(T)} \det(sI - A) \det(T) = \det(sI - A)
    \end{align}
    Since the characteristic polynomials are identical, their roots (the eigenvalues) are identical.
    \item \textbf{Transfer Function}:
    \begin{align}
        \tilde{C}(sI - \tilde{A})^{-1}\tilde{B} + \tilde{D} &= (CT)(sI - T^{-1}AT)^{-1}(T^{-1}B) + D \nonumber \\
        &= (CT)(T^{-1}(sI - A)T)^{-1}(T^{-1}B) + D \nonumber \\
        &= (CT)(T^{-1}(sI - A)^{-1}T)(T^{-1}B) + D \nonumber \\
        &= C (T T^{-1}) (sI - A)^{-1} (T T^{-1}) B + D \nonumber \\
        &= C (sI - A)^{-1} B + D
    \end{align}
\end{enumerate}
\end{proof}

\section{Transformation to Diagonal and Real Modal Canonical Forms}

\subsection{Diagonalization}
If a matrix $A \in \R^{n \times n}$ has $n$ linearly independent eigenvectors $v_1, \dots, v_n$, we can construct the non-singular modal matrix $T = [v_1 \dots v_n]$. The similarity transformation yields a diagonal dynamic matrix:
\begin{equation}
    \tilde{A} = T^{-1} A T = \Lambda = \diag(\lambda_1, \lambda_2, \dots, \lambda_n)
\end{equation}
In this coordinate system, the state variables are completely uncoupled, and the state trajectories evolve independently:
\begin{equation}
    \tilde{x}_i(t) = e^{\lambda_i t} \tilde{x}_i(0)
\end{equation}

\subsection{Real Modal Form for Complex Conjugate Eigenvalues}
If $A$ has complex conjugate eigenvalues $\lambda_{1,2} = \alpha \pm j \beta$, the corresponding eigenvectors $v_{1,2} = u \pm j w$ are also complex conjugates. To avoid working with complex-valued state-space matrices, we can use the transformation matrix:
\begin{equation}
    T = \begin{bmatrix} \text{Re}(v_1) & \text{Im}(v_1) \end{bmatrix} = \begin{bmatrix} u & w \end{bmatrix}
\end{equation}
The resulting similar matrix is in the real Jordan (modal) form:
\begin{equation}
    \tilde{A} = T^{-1} A T = \begin{bmatrix} \alpha & -\beta \\ \beta & \alpha \end{bmatrix}
\end{equation}
The matrix exponential for this block can be computed analytically, yielding a combined exponential decay and sinusoidal oscillation:
\begin{equation}
    \exp\left( \begin{bmatrix} \alpha & -\beta \\ \beta & \alpha \end{bmatrix} t \right) = e^{\alpha t} \begin{bmatrix} \cos(\beta t) & -\sin(\beta t) \\ \sin(\beta t) & \cos(\beta t) \end{bmatrix}
\end{equation}

\section{The Cayley-Hamilton Theorem \& Briot-Ruffini division}
The Cayley-Hamilton Theorem provides a highly systematic method for computing matrix functions, such as the matrix exponential, without relying on infinite series.

\begin{mytheorem}{Cayley-Hamilton Theorem}
Every square matrix $A \in \R^{n \times n}$ satisfies its own characteristic equation. That is, if $\Delta(\lambda) = \det(\lambda I - A) = \lambda^n + a_{n-1} \lambda^{n-1} + \dots + a_0$, then:
\begin{equation}
    \Delta(A) = A^n + a_{n-1} A^{n-1} + \dots + a_0 I = 0
\end{equation}
\end{mytheorem}

\subsection{Briot-Ruffini Division for Matrix Functions}
We can express any analytic matrix function $f(A)$ (such as $e^{At}$) using a polynomial of degree at most $n-1$. Let $a(\lambda) = e^{\lambda t}$ be our target function, and divide it by the characteristic polynomial $\Delta(\lambda)$ using polynomial division (Briot-Ruffini):
\begin{equation}
    e^{\lambda t} = \Delta(\lambda) q(\lambda, t) + r(\lambda, t)
\end{equation}
where the remainder $r(\lambda, t)$ has degree at most $n-1$:
\begin{equation}
    r(\lambda, t) = \sum_{i=0}^{n-1} \rho_i(t) \lambda^i = \rho_0(t) + \rho_1(t) \lambda + \dots + \rho_{n-1}(t) \lambda^{n-1}
\end{equation}
Evaluating this relation at the eigenvalues $\lambda_j$ of $A$ (which are the roots of $\Delta(\lambda) = 0$) gives a system of algebraic equations to solve for the coefficients $\rho_i(t)$:
\begin{equation}
    e^{\lambda_j t} = r(\lambda_j, t) = \sum_{i=0}^{n-1} \rho_i(t) \lambda_j^i, \quad j=1,\dots,n
\end{equation}
By the Cayley-Hamilton Theorem, evaluating the polynomial division at the matrix $A$ yields $\Delta(A) = 0$, leaving:
\begin{equation}
    e^{At} = r(A, t) = \rho_0(t) I + \rho_1(t) A + \dots + \rho_{n-1}(t) A^{n-1}
\end{equation}
This formula reduces the computation of the matrix exponential of an $n \times n$ matrix to the calculation of $n$ scalar coefficients and a finite sum of matrix powers!

\section{Jordan Canonical Forms \& Generalized Eigenvectors}
When a matrix $A$ has repeated eigenvalues, it may not possess $n$ linearly independent eigenvectors. In such cases, the matrix is said to be \textbf{defective} and cannot be diagonalized. The closest representation to a diagonal structure is the \textbf{Jordan Canonical Form}.

\subsection{Algebraic and Geometric Multiplicity}
\begin{itemize}
    \item \textbf{Algebraic Multiplicity} ($n_i$): The number of times the eigenvalue $\lambda_i$ appears as a root of the characteristic characteristic polynomial $\det(sI - A) = 0$.
    \item \textbf{Geometric Multiplicity} ($q_i$): The number of linearly independent eigenvectors associated with $\lambda_i$. It is the dimension of the null space (kernel) of $(A - \lambda_i I)$:
    \begin{equation}
        q_i = \dim \{ \ker(A - \lambda_i I) \} = n - \text{rank}(A - \lambda_i I)
    \end{equation}
\end{itemize}
For any eigenvalue, $1 \le q_i \le n_i$. If $q_i < n_i$ for any eigenvalue, the matrix is defective.

\subsection{The Jordan Block and Generalized Eigenvectors}
A defective matrix is transformed into a block-diagonal matrix where each diagonal element is a **Jordan Block**.
\begin{mydefinition}{Jordan Block}
A Jordan block of order $k$ associated with eigenvalue $\lambda$ is a $k \times k$ upper-triangular matrix of the form:
\begin{equation}
    J_k(\lambda) = \begin{bmatrix}
        \lambda & 1 & 0 & \cdots & 0 \\
        0 & \lambda & 1 & \cdots & 0 \\
        \vdots & \vdots & \ddots & \ddots & \vdots \\
        0 & 0 & 0 & \lambda & 1 \\
        0 & 0 & 0 & 0 & \lambda
    \end{bmatrix}
\end{equation}
\end{mydefinition}
To construct the similarity transformation matrix $T$ for a Jordan Block, we must calculate **generalized eigenvectors**.
For a Jordan block of size $k$ with eigenvalue $\lambda$:
\begin{enumerate}
    \item The first vector $v_1$ is a standard eigenvector satisfying:
    \begin{equation}
        (A - \lambda I) v_1 = 0
    \end{equation}
    \item The subsequent generalized eigenvectors $v_2, \dots, v_k$ are calculated recursively, forming a **Jordan Chain**:
    \begin{align}
        (A - \lambda I) v_2 &= v_1 \\
        (A - \lambda I) v_3 &= v_2 \\
        \vdots \nonumber \\
        (A - \lambda I) v_k &= v_{k-1}
    \end{align}
\end{enumerate}

\subsection{Matrix Exponential of a Jordan Block}
A Jordan block $J_k(\lambda)$ can be split into a diagonal matrix and a nilpotent matrix:
\begin{equation}
    J_k(\lambda) = \lambda I_k + N_k, \quad \text{where } N_k = \begin{bmatrix}
        0 & 1 & 0 & \cdots & 0 \\
        0 & 0 & 1 & \cdots & 0 \\
        \vdots & \vdots & \ddots & \ddots & \vdots \\
        0 & 0 & 0 & 0 & 1 \\
        0 & 0 & 0 & 0 & 0
    \end{bmatrix}
\end{equation}
Since $\lambda I_k$ and $N_k$ commute ($\lambda I_k N_k = N_k \lambda I_k$), we can split the matrix exponential:
\begin{equation}
    e^{J_k(\lambda) t} = e^{\lambda I_k t} e^{N_k t} = e^{\lambda t} e^{N_k t}
\end{equation}
Because $N_k$ is a nilpotent matrix of order $k$ (meaning $N_k^k = 0$), the infinite series for $e^{N_k t}$ terminates after $k$ terms:
\begin{equation}
    e^{N_k t} = I + N_k t + \frac{1}{2!} N_k^2 t^2 + \dots + \frac{1}{(k-1)!} N_k^{k-1} t^{k-1}
\end{equation}
Evaluating this sum yields the analytical matrix exponential for a Jordan Block:
\begin{equation}
    e^{J_k(\lambda) t} = e^{\lambda t} \begin{bmatrix}
        1 & t & \frac{t^2}{2!} & \cdots & \frac{t^{k-1}}{(k-1)!} \\
        0 & 1 & t & \cdots & \frac{t^{k-2}}{(k-2)!} \\
        \vdots & \vdots & \ddots & \ddots & \vdots \\
        0 & 0 & 0 & 1 & t \\
        0 & 0 & 0 & 0 & 1
    \end{bmatrix}
\end{equation}
This derivation reveals that repeated eigenvalues with geometric multiplicity 1 introduce algebraic polynomial terms $t, t^2, \dots$ multiplying the natural exponential modes $e^{\lambda t}$ in the time domain, which has significant implications for system stability.
